\subsection{非紧区间上积分对参数的导数}

设集 V 与 II, p.~74 的命题 5 中的含义相同，现在假设 I 是 R 中的任意区间，f 是 I × V 到 E 中的连续映射；若积分 \(\mathbf{g}(\alpha) = \int_{\mathrm{I}} \mathbf{f}(t, \alpha) dt\) 对所有 \(\alpha \in V\) 存在且是 \(\alpha\) 的连续函数，则函数 g 未必有一个等于 \(\int_{\mathrm{I}} \mathbf{f}'_{\alpha}(t, \alpha_0) dt\) 的导数（在点 \(\alpha_0\) 处），即使 \(\mathbf{f}'_{\alpha}(x, \alpha)\)

在含于 I 的每个紧区间上一致收敛于 \(\mathbf{f}_{\alpha}^{\prime}(x,\alpha_{0})\)，且积分 \(\int_{\mathrm{I}}\mathbf{f}_{\alpha}^{\prime}(t,\alpha)dt\) 对所有 \(\alpha \in \mathrm{V}\) 存在（参见 II, p.~87, 习题 3）。

下面的命题给出了使公式 (12)（II, p.~74）仍然成立的一个充分条件：

命题 7. 设 I 是 \(\mathbf{R}\) 中的任意区间，\(\mathbf{f}\) 是 \(\mathrm{I} \times \mathrm{V}\) 上的连续函数。假设：

\(1^{\circ}\) 偏导数 \(\mathbf{f}_{\alpha}^{\prime}(x,\alpha)\) 对所有 \(x\in \mathrm{I}\) 与所有 \(\alpha \in \mathrm{V}\) 存在，且对所有 \(\alpha \in \mathrm{V}\)，映射 \(x\mapsto \mathbf{f}_{\alpha}^{\prime}(x,\alpha)\) 在 \(\mathrm{I}\) 上是规则的；

\(2^{\circ}\) 对所有 \(\alpha \in \mathrm{V}\)，\(\mathbf{f}_{\alpha}^{\prime}(x,\beta)\) 一致收敛于 \(\mathbf{f}_{\alpha}^{\prime}(x,\alpha)\)（在含于 I 的每个紧区间上，当 \(\beta\) 趋于 \(\alpha\) 时）；

\(3^{\circ}\) 积分 \(\int_{\mathrm{I}}\mathbf{f}_{\alpha}^{\prime}(t,\alpha)dt\) 在 \(\mathrm{V}\) 上一致收敛；

\(4^{\circ}\) 积分 \(\int_{\mathrm{I}}\mathbf{f}(t,\alpha_0)dt\) 收敛。

在这些条件之下，积分 \(\mathbf{g}(\alpha)=\int_{\mathrm{I}}\mathbf{f}(t,\alpha)dt\) 在 V 上一致收敛，且函数 g 在 V 的每一点处有一个由下式给出的导数

\[
\mathbf {g} ^ {\prime} (\alpha) = \int_ {\mathrm{I}} \mathbf {f} _ {\alpha} ^ {\prime} (t, \alpha) d t.\tag{14}
\]

\(\int_{\mathrm{I}} \mathbf{f}_{\alpha}'(t, \alpha) dt\) 在 V 上的一致收敛意味着：函数 \(\alpha \mapsto \int_{\mathrm{J}} \mathbf{f}_{\alpha}'(t, \alpha) dt\) 在 V 上关于滤子 \(\Phi\)——即有向集 \(\Re(\mathrm{I})\)（由含于 I 的紧区间 J 组成）的截口滤子——一致收敛。令 \(\mathbf{u}_{\mathrm{J}}(\alpha) = \int_{\mathrm{J}} \mathbf{f}(t, \alpha) dt\)；这些假设表明，一方面 \(\mathbf{u}_{\mathrm{J}}(\alpha_0)\) 关于 \(\Phi\) 有一个极限，另一方面，由 II, p.~74 的命题 5，有 \(\mathbf{u}_{\mathrm{J}}'(\alpha) = \int_{\mathrm{J}} \mathbf{f}_{\alpha}'(t, \alpha) dt\)（对一切 \(\alpha \in \mathrm{V}\)）。因此我们可以把 II, p.~52 的定理 1 应用于函数 \(\mathbf{u}_{\mathrm{J}}\)，这里指标集的角色由 \(\Re(\mathrm{I})\) 扮演，而这个集合上的滤子的角色由滤子 \(\Phi\) 扮演；命题随之立得。

注记。1) 命题 7 的条件 \(1^{\circ}\) 与 \(2^{\circ}\) 更是满足的，当 \(\mathbf{f}_{\alpha}^{\prime}(x,\alpha)\) 是 \((x,\alpha)\) 在 \(\mathrm{I} \times \mathrm{V}\) 上的连续函数时。

\begin{enumerate}
\def\labelenumi{\arabic{enumi})}
\setcounter{enumi}{1}
\tightlist
\item
  在积分 \(\int_{a(\alpha)}^{b(\alpha)}\mathbf{f}(t,\alpha)dt\) 中，当区间的端点是参数的有限函数时，把这个积分作为 \(\alpha\) 的函数来研究可以化为对 {[}0, 1{]} 上积分的研究；事实上，通过变量替换 \(t = a(\alpha)(1 - u) + b(\alpha)u\)，有
\end{enumerate}

\[
\int_ {a (\alpha)} ^ {b (\alpha)} \mathbf {f} (t, \alpha) d t = \int_ {0} ^ {1} \mathbf {f} (a (\alpha) (1 - u) + b (\alpha) u, \alpha) (b (\alpha) - a (\alpha)) d u.
\]

